\documentclass[14pt]{extarticle}
\usepackage{natbib}

\usepackage{xcolor}
\usepackage{amsmath}
\newcommand{\tuple}[1]{ \langle #1 \rangle }
%\usepackage{automata}
\usepackage{times}
\usepackage{ltablex}
\usepackage{tasks}
\usepackage{threeparttable}
\usepackage{booktabs}
\usepackage{xcolor} % For custom colors
\usepackage{tikz} % For styling enumerate numbers
\usepackage{tcolorbox} % For colored box styling
\usepackage{amsmath, amsfonts, amssymb, amsthm} % Math related
\usepackage{natbib}
\usepackage{fontspec}
\usepackage{luatexja}
\usepackage[mathscr]{euscript}

%%%%%% Template
\usepackage{hyperref}
\hypersetup{colorlinks=true,allcolors=blue}

\usepackage{vmargin}
\setpapersize{USletter}
\setmarginsrb{1.0in}{1.0in}{1.0in}{0.6in}{0pt}{0pt}{0pt}{0.4in}

% HOW TO USE THE ABOVE:
%\setmarginsrb{leftmargin}{topmargin}{rightmargin}{bottommargin}{headheight}{headsep}{footheight}{footskip}
%\raggedbottom
% paragraphs indent & skip:
\parindent  0.3cm
\parskip    -0.01cm

\usepackage{tikz}
\usetikzlibrary{backgrounds}

% hyphenation:
% \hyphenpenalty=10000 % no hyphen
% \exhyphenpenalty=10000 % no hyphen
\sloppy

% notes-style paragraph spacing and indentation:
\usepackage{parskip}
\setlength{\parindent}{0cm}

% let derivations break across pages
\allowdisplaybreaks

\newcommand{\orange}[1]{\textcolor{orange}{#1}}
\newcommand{\blue}[1]{\textcolor{blue}{#1}}
\newcommand{\red}[1]{\textcolor{red}{#1}}
\newcommand{\freq}[1]{{\bf \sf F}(#1)}
\newcommand{\datafreq}[2]{{{\bf \sf F}_{#1}(#2)}}

\def\qqquad{\quad\qquad}
\def\qqqquad{\qquad\qquad}

%%%%%%%%%%%%%%%%%%%%%%%%%%%%%%%%%%%%%%%%%%%%%%%%%%%%%%%%%%%%%%%%%%%%%%%%%%%%%%%%
%%%%%%%%%%%%%%%%%%%%%%%%%%%%%%%%%%%%%%%%%%%%%%%%%%%%%%%%%%%%%%%%%%%%%%%%%%%%%%%%

% fill-in-blank question style, found in https://tex.stackexchange.com/a/505089

\usepackage{ifthen}
\usepackage{tocloft}
\usepackage{exercise}
% \usepackage{xcolor}

% Set the Show Answers Boolean
\newboolean{showAns}
\setboolean{showAns}{false}
\newcommand{\showAns}{\setboolean{showAns}{true}}

% The length of the Answer line
\newlength{\answerlength}
\newcommand{\anslen}[1]{\settowidth{\answerlength}{#1}}

% ans command that indicates space for an answer or shows the answer in red
\newcommand{\ans}[1]{\settowidth{\answerlength}{\hspace{2ex}#1\hspace{2ex}}%
    \ifthenelse{\boolean{showAns}}%
        {\textcolor{red}{\underline{\hspace{2ex}#1\hspace{2ex}}}}%
        {\underline{\hspace{\answerlength}}}}%

\newcommand{\details}[1]{\settowidth{\answerlength}{#1}%
    \ifthenelse{\boolean{showAns}}%
        {\begin{tcolorbox}[colback=blue!5!white,colframe=blue!75!black,title=Solution] #1 \end{tcolorbox}}%
        {}}%

% Formatting how multiple choices Questions are formated.
\settasks{label=(\Alph*), label-width=30pt}


% Some commands for the Exercise Question package
\renewcommand{\QuestionNB}{\Large\protect\textcircled{\small\bfseries\arabic{Question}}\ }
\renewcommand{\ExerciseHeader}{} %no header
\renewcommand{\QuestionBefore}{3ex} %Space above each Q
\setlength{\QuestionIndent}{8pt} % Indent after Q number


% To create the list of answers with tocloft...
\newcommand{\listanswername}{Answers}
\newlistof[Question]{answer}{Answers}{\listanswername}

% Creates a TOC for Answers
\newcounter{prevQ}
\newcommand{\answer}[1]{\refstepcounter{answer}%
\ans{#1}%
\ifnum\theQuestion=\theprevQ%
        \addcontentsline{Answers}{answer}{\protect\numberline{}#1}% don't include the Q number
        \else%
        \addcontentsline{Answers}{answer}{\protect\numberline{\theQuestion}#1}%
        \setcounter{prevQ}{\value{Question}}%
        \fi%
        }%

% \hyphenpenalty=10000 % no hyphen
% \exhyphenpenalty=10000 % no hyphen
\sloppy              % hyphen

\newcommand{\HRule}{\rule{\linewidth}{0.5mm}}
\newcommand{\Hrule}{\rule{\linewidth}{0.3mm}}

%tocloft formatting listofanswers
\renewcommand{\cftAnswerstitlefont}{\bfseries\large}
\renewcommand{\cftanswerdotsep}{\cftnodots}
\cftpagenumbersoff{answer}
\addtolength{\cftanswernumwidth}{10pt}

\makeatletter% since there's an at-sign (@) in the command name
\renewcommand{\@maketitle}{%
  \parindent=0pt% don't indent paragraphs in the title block
  \centering
  {\Large \bfseries\textsc{\@title}} \\
  \vspace{5pt}
  {\large \textit{\@author}} \\
  \HRule \\
  \vspace{1em}
}
\makeatother% resets the meaning of the at-sign (@)


% --- Fonts (robust fallback) ---
\IfFontExistsTF{Crimson Pro Light}{
  \setmainfont{Crimson Pro Light}[
    ItalicFont={* Italic},
    BoldFont={Crimson Pro Medium},
    BoldItalicFont={Crimson Pro Medium Italic}]
  \setsansfont{Crimson Pro Light}[
    ItalicFont={* Italic},
    BoldFont={Crimson Pro Medium},
    BoldItalicFont={Crimson Pro Medium Italic}]
}{
  \setmainfont{HaranoAjiMincho}
  \setsansfont{HaranoAjiGothic}
}
\title{Final Exam}
\author{Macroeconomics II \\ Hui-Jun Chen}


% ------------------------------------------------------------
% Explanations for each option (only printed when \showAns is enabled)
% Usage: \ExplainChoices{A}{B}{C}{D}
% ------------------------------------------------------------
\newcommand{\ExplainChoices}[4]{%
\ifthenelse{\boolean{showAns}}{%
\vspace{0.3em}\noindent\textcolor{red}{\textbf{Explanations.}}
\begin{itemize}
\item[(A)] #1
\item[(B)] #2
\item[(C)] #3
\item[(D)] #4
\end{itemize}
}{}%
}
\begin{document}

\maketitle

% Uncomment to display answers + option-by-option explanations
% \showAns

\begin{Exercise}
\section*{Problem 1 (Questions 1--20): Expectations, NKPC, shocks}

\Question \textbf{Definition. If $p_t\equiv\log P_t$, then inflation is}
\answer{A}
\begin{tasks}(1)
\task $\pi_t\equiv p_t-p_{t-1}$
\task $\pi_t\equiv P_t-P_{t-1}$
\task $\pi_t\equiv \log P_t$
\task $\pi_t\equiv p_{t+1}-p_t$
\end{tasks}
\ExplainChoices{Correct.}{Wrong: depends on units; not the model definition.}{Wrong: that is the log price level.}{Wrong timing.}

\Question \textbf{Three inflation objects. Which pair must be distinguished in NK models?}
\answer{B}
\begin{tasks}(1)
\task $\pi_t$ and $i_t$
\task $\pi_t$ and $\mathbb{E}_t\pi_{t+1}$
\task $x_t$ and $u_t$
\task $\beta$ and $\kappa$
\end{tasks}
\ExplainChoices{Not the key expectations distinction.}{Correct.}{Different objects, but not the core expectations issue.}{Parameters.}

\Question \textbf{Adaptive expectations. Under adaptive expectations, we set}
\answer{B}
\begin{tasks}(1)
\task $\pi_t^e=\mathbb{E}_t\pi_{t+1}$
\task $\pi_t^e=\pi_{t-1}$
\task $\pi_t^e=\pi^*$ always
\task $\pi_t^e=0$ always
\end{tasks}
\ExplainChoices{Forward-looking closure.}{Correct.}{Anchored benchmark.}{Not standard.}

\Question \textbf{Forward-looking closure. A convenient NK choice of is}
\answer{B}
\begin{tasks}(1)
\task $\pi_t^e=\pi_{t-1}$
\task $\pi_t^e=\mathbb{E}_t\pi_{t+1}$
\task $\pi_t^e=\pi_t$
\task $\pi_t^e=i_t$
\end{tasks}
\ExplainChoices{Adaptive.}{Correct.}{Not an expectation.}{Wrong.}

\Question \textbf{NKPC. The New Keynesian Phillips Curve is}
\answer{A}
\begin{tasks}(1)
\task $\pi_t=\beta\,\mathbb{E}_t\pi_{t+1}+\kappa x_t+u_t$
\task $\pi_t=\pi_{t-1}+\kappa x_t+u_t$
\task $\pi_t=\kappa x_t+u_t$
\task $\pi_t=\beta\,\pi_{t-1}-\kappa x_t$
\end{tasks}
\ExplainChoices{Correct.}{Backward-looking PC.}{Missing expectations.}{Wrong sign/objects.}

\Question \textbf{Expectations matter even if slack is zero. If $x_t=0$ and $u_t=0$, NKPC implies}
\answer{B}
\begin{tasks}(1)
\task $\pi_t=0$ always
\task $\pi_t=\beta\,\mathbb{E}_t\pi_{t+1}$
\task $\pi_t=\pi_{t-1}$
\task $\pi_t=\kappa$
\end{tasks}
\ExplainChoices{Wrong unless expected inflation is zero.}{Correct.}{Adaptive closure.}{Wrong.}

\Question \textbf{Meaning of $x_t$. In this course, the output gap is}
\answer{B}
\begin{tasks}(1)
\task $x_t\equiv y_t+y_t^{FE}$
\task $x_t\equiv y_t-y_t^{FE}$
\task $x_t\equiv \pi_t-\pi^*$
\task $x_t\equiv i_t-\pi_t$
\end{tasks}
\ExplainChoices{Wrong.}{Correct.}{Inflation gap.}{Real-rate object.}

\Question \textbf{Interpretation of $\beta$ in NKPC. A larger $\beta$ means}
\answer{A}
\begin{tasks}(1)
\task more forward-looking inflation (greater weight on expected future inflation)
\task steeper NKPC (larger response to $x_t$)
\task tougher monetary policy
\task larger cost-push shocks
\end{tasks}
\ExplainChoices{Correct.}{That is $\kappa$.}{That is $\phi_\pi$.}{That is $u_t$.}

\Question \textbf{Interpretation of $\kappa$ in NKPC. A larger $\kappa$ means}
\answer{B}
\begin{tasks}(1)
\task more weight on expected inflation
\task inflation responds more to slack $x_t$
\task the central bank reacts more to inflation
\task the natural real rate rises
\end{tasks}
\ExplainChoices{That is $\beta$.}{Correct.}{That is $\phi_\pi$.}{Not implied.}

\Question \textbf{Cost-push shock. In NKPC, $u_t$ is best interpreted as}
\answer{B}
\begin{tasks}(1)
\task a demand shock moving the IS curve
\task a markup/energy/bottleneck shock shifting inflation for given $x_t$
\task the inflation target
\task the nominal interest rate
\end{tasks}
\ExplainChoices{Wrong block.}{Correct.}{Target is $\pi^*$.}{Wrong.}

\Question \textbf{Calvo pricing. Each period, a fraction of firms that can reset prices is}
\answer{B}
\begin{tasks}(1)
\task $\theta$
\task $1-\theta$
\task $\beta$
\task $\kappa$
\end{tasks}
\ExplainChoices{Wrong: $\theta$ keep price fixed.}{Correct.}{Wrong.}{Wrong.}

\Question \textbf{Calvo intuition. A firm that resets at time $t$ chooses $P_t^\#$ to maximize profits while}
\answer{A}
\begin{tasks}(1)
\task the price remains fixed with some probability
\task money supply is constant
\task inflation is exactly zero
\task the output gap is zero
\end{tasks}
\ExplainChoices{Correct.}{Not the condition.}{Not the objective.}{Not required.}

\Question \textbf{Shift in NKPC. Holding $x_t$ fixed, an increase in $\mathbb{E}_t\pi_{t+1}$}
\answer{B}
\begin{tasks}(1)
\task lowers $\pi_t$ one-for-one
\task raises $\pi_t$ by $\beta$ times the change
\task does not affect $\pi_t$
\task changes $\kappa$
\end{tasks}
\ExplainChoices{Wrong sign.}{Correct.}{Wrong.}{Wrong.}

\Question \textbf{Demand shock comovement. A positive demand shock (raising $x_t$) tends to generate}
\answer{A}
\begin{tasks}(1)
\task $x_t\uparrow$ and $\pi_t\uparrow$
\task $x_t\downarrow$ and $\pi_t\uparrow$
\task $x_t\uparrow$ and $\pi_t\downarrow$
\task $x_t\downarrow$ and $\pi_t\downarrow$
\end{tasks}
\ExplainChoices{Correct.}{Cost-push pattern.}{Wrong.}{Contractionary demand.}

\Question \textbf{Cost-push comovement. A positive cost-push shock $u_t>0$ (with stabilization) tends to generate}
\answer{B}
\begin{tasks}(1)
\task $x_t\uparrow$ and $\pi_t\uparrow$
\task $x_t\downarrow$ and $\pi_t\uparrow$
\task $x_t\uparrow$ and $\pi_t\downarrow$
\task $x_t\downarrow$ and $\pi_t\downarrow$
\end{tasks}
\ExplainChoices{Not typical.}{Correct.}{Wrong.}{Wrong.}

\Question \textbf{Calculation (NKPC). Let $\beta=0.98$, $\mathbb{E}_t\pi_{t+1}=3\%$, $\kappa=0.5$, $x_t=-2\%$, and $u_t=0$. Then $\pi_t$ is}
\answer{A}
\begin{tasks}(1)
\task $1.94\%$
\task $2.94\%$
\task $3.94\%$
\task $0.50\%$
\end{tasks}
\ExplainChoices{Correct: $0.98\cdot3-1=1.94$.}{Wrong: ignores $x_t$.}{Wrong sign.}{Wrong.}

\Question \textbf{Calculation (solve for slack). Suppose $\pi_t=3\%$, $\beta=1$, $\mathbb{E}_t\pi_{t+1}=2\%$, $\kappa=0.5$, and $u_t=0$. Then $x_t$ equals}
\answer{A}
\begin{tasks}(1)
\task $2\%$
\task $1\%$
\task $0\%$
\task $-1\%$
\end{tasks}
\ExplainChoices{Correct: $3-2=0.5x\Rightarrow x=2$.}{Wrong.}{Wrong.}{Wrong.}

\Question \textbf{Stabilizing a one-time cost-push shock. If $\mathbb{E}_t\pi_{t+1}=\pi^*$ and the central bank wants $\pi_t=\pi^*$ with $u_t=1\%$ and $\kappa=0.25$, then it must set}
\answer{A}
\begin{tasks}(1)
\task $x_t=-4\%$
\task $x_t=-1\%$
\task $x_t=+4\%$
\task $x_t=0\%$
\end{tasks}
\ExplainChoices{Correct: $x=-u/\kappa=-4$.}{Wrong magnitude.}{Wrong sign.}{Wrong.}

\Question \textbf{Anchored vs drifting. Expectations are drifting if}
\answer{B}
\begin{tasks}(1)
\task after a one-time shock, $\mathbb{E}_t\pi_{t+1}$ returns to $\pi^*$
\task after a one-time shock, $\mathbb{E}_t\pi_{t+1}$ stays above $\pi^*$
\task $x_t=0$ always
\task $u_t=0$ always
\end{tasks}
\ExplainChoices{Anchored.}{Correct.}{Not definition.}{Not definition.}

\Question \textbf{Key contrast. Relative to adaptive expectations, forward-looking expectations make inflation more sensitive to}
\answer{B}
\begin{tasks}(1)
\task last period inflation mechanically
\task beliefs about future policy (credibility) today
\task money demand only
\task the money supply only
\end{tasks}
\ExplainChoices{Adaptive feature.}{Correct.}{Wrong.}{Wrong.}

\vspace{0.6em}\Hrule\vspace{0.8em}
\section*{Problem 2 (Questions 21--40): NK IS, Taylor rule, determinacy, ZLB}

\Question \textbf{Fisher (ex-ante real rate). The expected real policy rate is approximately}
\answer{B}
\begin{tasks}(1)
\task $i_t+\mathbb{E}_t\pi_{t+1}$
\task $i_t-\mathbb{E}_t\pi_{t+1}$
\task $\pi_t-i_t$
\task $i_t-\pi_t$
\end{tasks}
\ExplainChoices{Wrong sign.}{Correct.}{Wrong sign.}{Uses current inflation.}

\Question \textbf{NK IS (Euler) equation. A standard gap-form NK IS is}
\answer{A}
\begin{tasks}(1)
\task $x_t=\mathbb{E}_t x_{t+1}-\frac{1}{\sigma}\left(i_t-\mathbb{E}_t\pi_{t+1}-r_t^n\right)$
\task $x_t=\mathbb{E}_t x_{t+1}+\frac{1}{\sigma}\left(i_t-\mathbb{E}_t\pi_{t+1}-r_t^n\right)$
\task $x_t=-\sigma\left(i_t-\mathbb{E}_t\pi_{t+1}-r_t^n\right)$
\task $x_t=\kappa\left(i_t-\mathbb{E}_t\pi_{t+1}\right)$
\end{tasks}
\ExplainChoices{Correct.}{Wrong sign.}{Wrong scaling.}{Wrong.}

\Question \textbf{Interpretation of $\sigma$. A larger $\sigma$ implies the output gap responds}
\answer{B}
\begin{tasks}(1)
\task more strongly to the real-rate gap
\task less strongly to the real-rate gap
\task only to $u_t$
\task only to $\kappa$
\end{tasks}
\ExplainChoices{Slope is $1/\sigma$.}{Correct.}{Wrong.}{Wrong.}

\Question \textbf{Natural real rate. $r_t^n$ is best interpreted as}
\answer{B}
\begin{tasks}(1)
\task the nominal policy rate
\task the flexible-price real rate that would make $x_t=0$
\task the inflation target
\task the cost-push shock
\end{tasks}
\ExplainChoices{Wrong.}{Correct.}{Wrong.}{Wrong.}

\Question \textbf{Calculation (NK IS). Suppose $\mathbb{E}_t x_{t+1}=0$, $\sigma=2$, $i_t=4\%$, $\mathbb{E}_t\pi_{t+1}=2\%$, and $r_t^n=1\%$. Then $x_t$ is}
\answer{A}
\begin{tasks}(1)
\task $-0.5\%$
\task $-1\%$
\task $+1\%$
\task $0\%$
\end{tasks}
\ExplainChoices{Correct.}{Forgot $1/\sigma$.}{Wrong sign.}{Wrong.}

\Question \textbf{Calculation (real rate). If $i_t=3\%$ and $\mathbb{E}_t\pi_{t+1}=1.5\%$, then the expected real rate is}
\answer{B}
\begin{tasks}(1)
\task $4.5\%$
\task $1.5\%$
\task $-1.5\%$
\task $3\%$
\end{tasks}
\ExplainChoices{Adds instead of subtracts.}{Correct.}{Wrong sign.}{Nominal.}

\Question \textbf{Forward guidance channel. Holding $i_t$ fixed, a higher $\mathbb{E}_t\pi_{t+1}$ implies}
\answer{B}
\begin{tasks}(1)
\task higher real rate and lower $x_t$
\task lower real rate and higher $x_t$
\task no change in $x_t$
\task $x_t$ must fall because expectations are higher
\end{tasks}
\ExplainChoices{Wrong sign.}{Correct.}{Wrong.}{Wrong.}

\Question \textbf{Taylor rule (general form). A standard rule in the NK core is}
\answer{A}
\begin{tasks}(1)
\task $i_t=\bar{i}+\phi_\pi(\pi_t-\pi^*)+\phi_x x_t$
\task $i_t=\pi_t-\phi_\pi x_t$
\task $i_t=r_t^n$ always
\task $i_t=\mathbb{E}_t\pi_{t+1}$
\end{tasks}
\ExplainChoices{Correct.}{Wrong.}{Wrong.}{Wrong.}

\Question \textbf{Taylor principle in NK. A standard sufficient condition for determinacy is}
\answer{B}
\begin{tasks}(1)
\task $\phi_\pi<1$
\task $\phi_\pi>1$
\task $\phi_x<0$
\task $\kappa=0$
\end{tasks}
\ExplainChoices{Wrong.}{Correct.}{Not the standard condition.}{Not sufficient.}

\Question \textbf{Policy experiment sequence. After a one-time cost-push shock $u_t>0$, the NK logic is}
\answer{A}
\begin{tasks}(1)
\task NKPC pushes $\pi_t\uparrow$; rule raises $i_t$; NK IS yields $x_t\downarrow$; NKPC then helps bring $\pi_t$ down
\task NKPC pushes $\pi_t\uparrow$; rule lowers $i_t$; $x_t\uparrow$; inflation falls immediately
\task rule does nothing; inflation returns automatically
\task NK IS implies $x_t\uparrow$ when $i_t$ rises
\end{tasks}
\ExplainChoices{Correct.}{Wrong.}{Wrong.}{Wrong.}

\Question \textbf{ZLB constraint is}
\answer{B}
\begin{tasks}(1)
\task $r_t\ge 0$
\task $i_t\ge 0$
\task $\pi_t\ge 0$
\task $x_t\ge 0$
\end{tasks}
\ExplainChoices{Wrong.}{Correct.}{Wrong.}{Wrong.}

\Question \textbf{ZLB calculation. A rule implies $i_t=-0.5\%$, but ZLB binds so $i_t=0$. If $\mathbb{E}_t\pi_{t+1}=1\%$, then the expected real rate is}
\answer{B}
\begin{tasks}(1)
\task $-1.5\%$
\task $-1\%$
\task $0\%$
\task $+1\%$
\end{tasks}
\ExplainChoices{Wrong: use $i=0$.}{Correct: $0-1=-1$.}{Wrong.}{Wrong sign.}

\Question \textbf{ZLB amplification. At the ZLB, a fall in expected inflation tends to}
\answer{B}
\begin{tasks}(1)
\task lower the real rate and raise demand
\task raise the real rate and lower demand
\task leave the real rate unchanged
\task eliminate the NKPC
\end{tasks}
\ExplainChoices{Wrong sign.}{Correct.}{Wrong.}{Wrong.}

\Question \textbf{Three-equation NK core consists of}
\answer{A}
\begin{tasks}(1)
\task IS (Euler), NKPC, Taylor rule
\task IS, LM, SRAS
\task AD, SRAS, LRAS
\task budget constraint, money demand, IS
\end{tasks}
\ExplainChoices{Correct.}{Wrong.}{Wrong.}{Wrong.}

\Question \textbf{Sacrifice ratio (simple). If expectations are anchored ($\mathbb{E}_t\pi_{t+1}=\pi^*$) and $u_t=0$, then reducing inflation by 1 pp requires}
\answer{A}
\begin{tasks}(1)
\task $x_t=-1/\kappa$ pp
\task $x_t=-\kappa$ pp
\task $x_t=+1/\kappa$ pp
\task $x_t=0$
\end{tasks}
\ExplainChoices{Correct.}{Wrong.}{Wrong sign.}{Wrong.}

\Question \textbf{Sacrifice ratio (numeric). If $\kappa=0.2$ and you want inflation 0.5 pp lower (anchored expectations, $u_t=0$), required $x_t$ is}
\answer{A}
\begin{tasks}(1)
\task $-2.5\%$
\task $-0.1\%$
\task $+2.5\%$
\task $-10\%$
\end{tasks}
\ExplainChoices{Correct.}{Wrong.}{Wrong sign.}{Wrong.}

\Question \textbf{2-period equilibrium (calc). Suppose $\mathbb{E}_0\pi_1=\pi^*=2\%$, $\mathbb{E}_0x_1=0$, $\beta=1$, $\kappa=0.5$, $u_0=1\%$. Policy: $i_0=r_0^n+\mathbb{E}_0\pi_1+\phi_\pi(\pi_0-\pi^*)+\phi_x x_0$ with $\sigma=1$, $\phi_\pi=2$, $\phi_x=1$. Then $\pi_0$ is closest to}
\answer{C}
\begin{tasks}(1)
\task $2.0\%$
\task $2.33\%$
\task $2.67\%$
\task $3.0\%$
\end{tasks}
\ExplainChoices{Wrong.}{Wrong.}{Correct: $\pi_0\approx2.67\%$.}{Too high.}

\Question \textbf{In the same 2-period equilibrium as previous question, $x_0$ is closest to}
\answer{C}
\begin{tasks}(1)
\task $0.0\%$
\task $-0.33\%$
\task $-0.67\%$
\task $-2.0\%$
\end{tasks}
\ExplainChoices{Wrong.}{Wrong.}{Correct: $x_0\approx-0.67\%$.}{Too large.}

\Question \textbf{Determinacy intuition. The Taylor principle helps anchor inflation because it ensures that if $\pi_t$ rises, policy}
\answer{B}
\begin{tasks}(1)
\task lowers the real rate, validating higher inflation
\task raises the real rate enough to create $x_t<0$, pushing inflation down via NKPC
\task keeps the real rate fixed at $r_t^n$ regardless of inflation
\task eliminates cost-push shocks
\end{tasks}
\ExplainChoices{Destabilizing.}{Correct.}{That is $\phi_\pi=1$; not anchoring.}{Policy cannot delete shocks.}

\Question \textbf{Old vs new Phillips curve. Relative to the accelerationist PC, the NKPC replaces}
\answer{A}
\begin{tasks}(1)
\task $\pi_{t-1}$ with $\mathbb{E}_t\pi_{t+1}$
\task $x_t$ with $y_t$
\task $u_t$ with $i_t$
\task $\kappa$ with $\beta$
\end{tasks}
\ExplainChoices{Correct.}{Wrong.}{Wrong.}{Wrong.}

\end{Exercise}
\end{document}
